\chapter{The Big Picture}
\section{The problem}
A library holds thousands of \emph{scanned} pages (almanacs, ledgers, tide tables). The scans are noisy, skewed, and contain no machine-readable text. A user asks a natural-language question. The system must return an answer that is
\begin{itemize}
 \item \textbf{accurate} (the right fact),
 \item \textbf{faithful / grounded} (supported by the documents, with citations),
 \item \textbf{honest} (says ``insufficient evidence'' instead of guessing),
 \item \textbf{safe} (resists instructions hidden in documents, leaks no personal data),
 \item \textbf{auditable} (every decision is logged).
\end{itemize}

\section{Vocabulary you must know}
\begin{longtable}{@{}p{3.2cm}p{11.3cm}@{}}
\toprule \textbf{Term} & \textbf{Meaning in this codebase} \\ \midrule \endhead
OCR / HTR & Optical / Handwritten Text Recognition: image of text $\to$ string. Baseline model: TrOCR. \\
Layout detection & Finding \emph{where} the text is (lines, headings, figures) before reading it. \\
Chunk & A piece of text (here a sliding window of words) that is embedded and retrieved. \\
Embedding & A vector (384 numbers) such that similar meaning $\Rightarrow$ similar vectors. \\
Dense retrieval & Find chunks whose embedding is closest to the question's embedding. \\
Cosine similarity & $\cos(a,b)=\frac{a\cdot b}{\|a\|\|b\|}$. With unit vectors it equals the dot product. Used as the \emph{score}. \\
Rerank & A slower, more accurate model (cross-encoder) reorders the top candidates. \\
RAG & Retrieval-Augmented Generation: an LLM answers using retrieved text as evidence. \\
Agent & A loop that \emph{observes}, \emph{decides} which action to take, \emph{acts}, and repeats. \\
Agentic (here) & The decision depends on an observed number: re-search when evidence is weak. \\
Grounded & Every claim is supported by a cited chunk. \\
Abstain & Refuse to answer because evidence is insufficient. \\
Hook / seam & A fixed named point where cross-cutting code runs. \\
PII & Personally identifiable information (email, phone, ID numbers). \\
Prompt injection & Text in a document that tries to give the model new instructions. \\
HITL & Human-in-the-loop: a person reviews/approves an action. \\
ECE & Expected Calibration Error: gap between confidence and accuracy. \\
RLVR & Reinforcement learning with \emph{verifiable} rewards (exact-match check, no learned judge). \\
\bottomrule
\end{longtable}

\section{Two phases}
\begin{figure}[h]\centering
\begin{tikzpicture}[node distance=5mm,
 box/.style={draw,rounded corners,fill=blue!6,minimum width=4.3cm,minimum height=7mm,font=\small,align=center},
 hk/.style={draw,rounded corners,fill=orange!15,minimum width=4.3cm,minimum height=6mm,font=\footnotesize\itshape,align=center},
 >={Stealth}]
\node[font=\bfseries] (ta) {PHASE A: build (once)};
\node[box,below=of ta] (l) {Stage 1: load + preprocess + enhance};
\node[box,below=of l] (lay) {Stage 2: layout detection};
\node[box,below=of lay] (ocr) {Stage 3: OCR};
\node[hk,below=of ocr] (h1) {hook: AFTER\_OCR (PII redaction)};
\node[box,below=of h1] (ch) {Stage 4: chunk $\to$ embed $\to$ index};
\node[box,below=of ch,fill=green!10] (idx) {data/index (faiss + chunks)};
\draw[->] (ta)--(l); \draw[->] (l)--(lay); \draw[->] (lay)--(ocr); \draw[->] (ocr)--(h1); \draw[->] (h1)--(ch); \draw[->] (ch)--(idx);
\node[font=\bfseries,right=3.2cm of ta] (tb) {PHASE B: answer (per question)};
\node[box,below=of tb] (q) {question};
\node[box,below=of q] (ret) {Stage 5: retrieve top-$k$ (score)};
\node[box,below=of ret,fill=yellow!15] (ag) {Stage 6: agent loop\\decide / act / observe};
\node[hk,below=of ag] (h2) {hook: BEFORE\_ANSWER (PII, grounding)};
\node[box,below=of h2] (syn) {synthesize cited answer (LLM)};
\node[box,below=of syn,fill=green!10] (ans) {Answer or ``insufficient evidence''};
\draw[->] (tb)--(q); \draw[->] (q)--(ret); \draw[->] (ret)--(ag); \draw[->] (ag)--(h2); \draw[->] (h2)--(syn); \draw[->] (syn)--(ans);
\draw[->,dashed] (idx.east) -- ++(1.1,0) |- (ret.west);
\draw[->,bend left=35] (ag.east) to node[right,font=\scriptsize]{weak? widen $k$} (ret.east);
\end{tikzpicture}
\caption{Build phase (left) writes the index; answer phase (right) reads it.}
\end{figure}

The code for this whole picture is just twenty lines, \c{pipeline.py}. Read it now; everything else in the book elaborates it.
\snip{pipeline}
\begin{why}
The stage order is \emph{fixed} on purpose: graders and CI rely on the same structure in every student project. \c{wiring.register\_all(cfg)} is called at the start of both functions so that the cross-cutting features are attached \emph{before} anything runs.
\end{why}
\begin{gotcha}
\c{answer()} rebuilds a \c{Retriever} on every call, and \c{register\_all} calls \c{hooks.clear()} first. If it did not clear, handlers would pile up and run several times per query. There is a test for this (\c{test\_seams\_clear\_between\_queries}).
\end{gotcha}

\section{Data flow, with the types}
\begin{center}
\begin{tabular}{@{}lll@{}}
\toprule Stage & Input & Output \\ \midrule
1 load\_pages & \c{cfg} & \c{list[Page]} \\
1 preprocess / enhance & \c{list[Page]} & \c{list[Page]} (new image paths) \\
2 layout.detect & \c{list[Page]} & \c{list[Region]} \\
3 ocr.transcribe & \c{list[Region]} & \c{list[Chunk]} (one per line) \\
4 chunk.split & \c{list[Chunk]} & \c{list[Chunk]} (windows) \\
4 embed.encode & \c{list[Chunk]} & \c{ndarray (n, 384)} \\
4 store.build & chunks, vectors & files on disk \\
5 retrieve & question \c{str} & \c{list[Chunk]} with \c{score} \\
6 agent.run & question \c{str} & \c{Answer} \\
\bottomrule
\end{tabular}
\end{center}
